\centerline{\bf \Huge 4 класс}

\problem{\textbf{1}} \textbf{Критерии:} \\
\textit{7 баллов} $-$ Найден любой верный пример.\\
\textit{0-5 баллов} $-$ Верный пример не найден, но имеются продвижения.\\
\textit{0 баллов} $-$ Решение отсутствует.\\
\textit{0 баллов} $-$ При обоснованных подозрениях на списывание участником из любого источника (Интернет, учебное пособие, ключи олимпиады и т.д.) член жюри имеет право указать на такие фрагменты с помощью восклицательных знаков на полях. \\

\problem{\textbf{2}} \textbf{Критерии:} \\
\textit{0-4 балла} $-$ Решение А.\\
\textit{0-1 балл} $-$ Составление текстового условия Б.\\
\textit{0-2 балла} $-$ Любое правильное решение Б.\\
\textit{0 баллов} $-$ Любой ответ без объяснений.\\
\textit{0 баллов} $-$ Решение отсутствует.\\
\textit{0 баллов} $-$ При обоснованных подозрениях на списывание участником из любого источника (Интернет, учебное пособие, ключи олимпиады и т.д.) член жюри имеет право указать на такие фрагменты с помощью восклицательных знаков на полях. \\


\problem{\textbf{3}} \textbf{Критерии суммируются:} \\
\textit{2 балла} $-$ Изначально было $90\cdot2=180$ кг рыбы.\\
\textit{2 балла} $-$ Судаков и лещей суммарно $180-46-30=104$ кг.\\
\textit{2 балла} $-$ Лещей $x$ кг, судаков $3x$ кг => $x+3x=104$ кг, $x=26$ кг.\\
\textit{1 балл} $-$ Судаков больше, чем лещей на $3x-x=2x$ кг => на $2\cdot26=52$ кг.\\
\textit{0 баллов} $-$ Любой ответ без объяснений.\\
\textit{0 баллов} $-$ Решение отсутствует.\\
\textit{0 баллов} $-$ При обоснованных подозрениях на списывание участником из любого источника (Интернет, учебное пособие, ключи олимпиады и т.д.) член жюри имеет право указать на такие фрагменты с помощью восклицательных знаков на полях. \\

\textbf{ИЛИ:}
\textit{7 баллов} $-$ Любое другое полное верное решение.\\


\problem{\textbf{4}} \textbf{Критерии:} \\
\textit{7 баллов} $-$ Полное верное решение.\\
\textit{6-7 баллов} $-$ Верное решение. Имеются небольшие недочёты, в целом не влияющие на решение.\\
\textit{5-6 баллов} $-$ Решение в целом верное. Однако оно содержит ошибки, либо пропущенные случаи, не влияющие на логику рассуждений.\\
\textit{4-3 балла} $-$ Верно рассмотрен один из двух (более сложный) существенных случаев. В том случае, когда решение делится на две равноценные части – решение одной из частей. 
\\
\textit{2-3 балла} $-$ Доказаны вспомогательные утверждения, помогающие в решении задачи.
\\
\textit{0-1 балл} $-$ Рассмотрены отдельные важные случаи при отсутствии решения (или при ошибочном решении).
\\
\textit{0 баллов} $-$ Решение неверное, продвижения отсутствуют.\\
\textit{0 баллов} $-$ Решение отсутствует.
\\
\textit{0 баллов} $-$ При обоснованных подозрениях на списывание участником из любого источника (Интернет, учебное пособие, ключи олимпиады и т.д.) член жюри имеет право указать на такие фрагменты с помощью восклицательных знаков на полях. \\

\problem{\textbf{5}} \textbf{Критерии:} \\
\textit{7 баллов} $-$ Найден любой верный пример.\\
\textit{0 баллов} $-$ Любое неверное решение.\\
\textit{0 баллов} $-$ Решение отсутствует.\\
\textit{0 баллов} $-$ При обоснованных подозрениях на списывание участником из любого источника (Интернет, учебное пособие, ключи олимпиады и т.д.) член жюри имеет право указать на такие фрагменты с помощью восклицательных знаков на полях. \\
